\documentclass[11pt,a4paper]{article}

\usepackage[a4paper,margin=2.4cm]{geometry}
\usepackage{fontspec}
\usepackage{xcolor}
\usepackage{booktabs}
\usepackage{array}
\usepackage{amsmath}
\usepackage{amssymb}
\usepackage{enumitem}
\usepackage{titlesec}
\usepackage{fancyhdr}
\usepackage{microtype}

% ---------- Thai fonts ----------
\setmainfont{TH Sarabun New}[Scale=1.24]
\setsansfont{Tahoma}
\newfontfamily\thaifont{TH Sarabun New}[Scale=1.24]

% ---------- Colours (match deck vocabulary) ----------
\definecolor{myblue}{RGB}{0,41,120}
\definecolor{barblue}{RGB}{0,90,170}
\definecolor{rulegray}{RGB}{120,120,120}

% ---------- Headings ----------
\titleformat{\section}{\Large\bfseries\color{myblue}}{\thesection}{0.8em}{}
\titleformat{\subsection}{\large\bfseries\color{myblue}}{\thesubsection}{0.8em}{}
\titlespacing{\section}{0pt}{14pt}{6pt}
\titlespacing{\subsection}{0pt}{10pt}{4pt}

% ---------- Header/footer ----------
\pagestyle{fancy}
\fancyhf{}
\fancyhead[L]{\small\color{rulegray} sg\_fault\_cycle160 — การวิเคราะห์การสลับโหมด}
\fancyhead[R]{\small\color{rulegray}\thepage}
\renewcommand{\headrulewidth}{0.4pt}

% ---------- Code-ish inline ----------
\newcommand{\code}[1]{\texttt{\small #1}}

\setlist{nosep,leftmargin=1.6em}

\title{\bfseries\color{myblue} ทำไม IBR 2--4 จึงไม่สลับเป็น GFM\\[2pt]
\large การวิเคราะห์ Figure 7 (Severity Index) และ Figure 8 (Converter Mode)\\
ของ scenario \code{sg\_fault\_cycle160}}
\author{}
\date{9 กันยายน 2569}

\begin{document}
\maketitle
\vspace{-1.2em}

\section*{คำถาม}
จาก Figure 7 ดัชนีความรุนแรง $S_i$ ของ IBR ทั้งสี่ตัวพุ่งทะลุเส้น threshold
$\Gamma_{\mathrm{on}} = 0.65$ หลายครั้ง แต่ Figure 8 กลับแสดงว่า IBR$_1$
(บัส 2, แถบน้ำเงิน) ถูกสลับเป็น GFM ตั้งแต่ตอน SG หลุด ($t = 20$ s) และมีเพียง
IBR$_3$ (บัส 6, แถบแดง) ที่ถูกเสริมเข้ามาช่วงสั้น ๆ ส่วน IBR$_2$ (บัส 3) กับ
IBR$_4$ (บัส 8) คงโหมด GFL ตลอด run คำถามคือ ทำไมตัวที่เหลือจึงไม่สลับ ทั้งที่ดัชนีของมัน ``ควร'' สลับ หรือมี
กลไกใดมาขัดขวาง

(ตารางตลอดบันทึกนี้ระบุชื่อคู่กับหมายเลขบัสกำกับเสมอ เพราะตรรกะ supervisor
ทำงานบนหมายเลขบัส [2, 3, 6, 8] ส่วน Figure 7/8 ใช้ชื่อ IBR$_1$--IBR$_4$
ตามลำดับนั้น)

\section*{คำตอบโดยย่อ}
\textbf{ไม่มีกลไกใดมา ``ขัด'' รายตัว} คำตอบอยู่ที่สถาปัตยกรรมการตัดสินใจของ
supervisor: $S_i$ ของแต่ละตัวไม่ได้เป็นตัวตัดสินว่า ``ตัวไหนสลับ'' — มันเป็นเพียง
\textbf{ทริกเกอร์ระดับระบบ} ที่ตอบคำถามว่า ``\emph{เมื่อไร} ระบบต้องการ
forming support เพิ่ม'' ส่วนคำถามว่า ``\emph{ตัวไหน} สลับ'' เป็นสิทธิ์ของ
\emph{authenticated selector table} ที่เลือก strict superset ที่เล็กที่สุด
ของชุด GFM ปัจจุบันเสมอ และก่อนจะถึงขั้นนั้น หลักฐานต้องผ่านเกตเวลาสองชั้น
(minimum-hold และ dwell) ซึ่งหลักฐานของ IBR$_2$--IBR$_4$ ไม่เคยยืนยันได้นานพอ
หลังจากการสลับครั้งแรก

ทุกตัวเลขในบันทึกนี้อ่านมาจาก cache ของ run จริง
(\code{output/diagnostics/ieee14\_scenario\_suite/sg\_fault\_cycle160.mat},
2{,}288 accepted samples) เทียบกับตรรกะใน
\code{+stability/ts\_simulate\_ibr\_hybrid.m} บรรทัด 846--903 และ
\code{+stability/select\_support\_augmentation\_candidate.m}

\section{ชั้นที่ 1: $S_i$ เป็นทริกเกอร์ระดับระบบ ไม่ใช่ตัวเลือกรายตัว}

ตรรกะ supervisor ทำงานดังนี้ (ไฟล์ \code{ts\_simulate\_ibr\_hybrid.m}):
\begin{itemize}
\item วัด $S_i$ ของทุกตัวที่ยัง online จากสองเทอมเท่านั้น:
      $S_i = \mathrm{sat}\bigl(\tfrac{1}{2}J_{V,i} + \tfrac{1}{2}J_f\bigr)$
      ($J_f$ ใช้ร่วมกันทั้งระบบผ่าน $f_{\mathrm{COI}}$)
\item เงื่อนไข ``เครียด'' คือ
      \code{up\_stress = any(}$S_i \geq \Gamma_{\mathrm{on}}$\code{)}
      \textbf{เฉพาะบนตัวที่ยังเป็น GFL} (บรรทัด 861--863) —
      เป็นค่าเดียวทั้งเกาะ ไม่มีการจัดอันดับรายตัว
\item เมื่อเงื่อนไขติดและครบ dwell แล้ว การเลือกตัวสลับถูกส่งต่อให้
      \code{select\_support\_augmentation\_candidate} ซึ่งค้นหาใน
      selector table (15 คอนฟิกูเรชันที่ตรวจสอบและรับรองไว้ล่วงหน้า)
      โดยมีกฎจัดอันดับ: \emph{จำนวน GFM ที่เพิ่มน้อยที่สุด} $\to$
      \emph{ส่วนเผื่อเสถียรภาพมากที่สุด} $\to$ เลข device น้อยที่สุด
      (ฟังก์ชัน \code{precedes}, บรรทัด 101--112)
\end{itemize}

ดังนั้นแม้ IBR$_2$--IBR$_4$ จะมี $S_i$ สูงกว่า IBR$_1$ ในหลายช่วง
ตารางก็ไม่เคยมองค่า $S_i$ เลย — มันมองเพียงว่าชุด GFM ปัจจุบันคืออะไร
และ superset ใด ``เล็กสุดที่ยังปลอดภัย''

\section{ชั้นที่ 2: เกตเวลาสองชั้นบล็อกช่วงแรกที่ทุกตัวพุ่งพร้อมกัน}

\subsection{หลักฐานจาก trajectory จริง}
ทันทีหลัง SG trip ที่ $t = 20$ s ดัชนีของทั้งสี่ตัวพุ่งถึง $S_i = 1.000$
ค้างนาน 248--539 ms (Table~\ref{tab:excursions}) หากไม่มีเกตเวลา
ทั้งสี่ตัว ``สมควร'' สลับพร้อมกันตามที่ตาเปล่าอ่านจาก Figure 7

\begin{table}[h]
\centering\small
\caption{ช่วงที่ $S_i > \Gamma_{\mathrm{on}} = 0.65$ ของแต่ละตัว (อ่านจาก cache)}
\label{tab:excursions}
\begin{tabular}{@{}clll@{}}
\toprule
ตัว (บัส) & รอบ $t = 20$ s & รอบ $t \approx 22$ s & รอบ fault $t = 60$ s \\
\midrule
IBR$_1$ (2) & 20.00--20.25 (248 ms) & 22.09--22.16 (68 ms), 22.62--22.88 (269 ms) & 60.00--60.15 (151 ms) \\
IBR$_2$ (3) & 20.00--20.25 (248 ms) & 21.99--22.09, 22.09--22.22, 22.48--23.13 (651 ms) & 60.00--60.15 (151 ms) \\
IBR$_3$ (6) & 20.00--20.39 (389 ms) & 22.09--22.78 (694 ms) & 60.00--60.15 (151 ms) \\
IBR$_4$ (8) & 20.00--20.54 (539 ms) & 22.09--22.19 (100 ms) & 60.00--60.15 (151 ms) \\
\bottomrule
\end{tabular}
\end{table}

\subsection{เกตที่ 1: minimum-hold}
บรรทัด 867--868 บังคับว่า supervisor ต้องรอ
$T_{\mathrm{minhold}} = 1.0$ s นับจากการ trip ($t_{\mathrm{trip}} = 20$ s)
จึงจะเริ่มสะสม dwell ได้ — ตลอดช่วง $t \in [20, 21)$ s ที่ทั้งสี่ตัวทะลุ
$\Gamma_{\mathrm{on}}$ พร้อมกัน ตัวจับเวลาทุกตัวถูกรีเซ็ตเป็นศูนย์อยู่ตลอด
(\code{support\_status = 'MINIMUM\_HOLD'})

\subsection{เกตที่ 2: dwell ต้องต่อเนื่อง}
เมื่อปลดล็อกที่ $t = 21.0$ s ปรากฏว่าค่า $S_i$ ของทุกตัว \emph{ตกลงมา
ต่ำกว่า} $\Gamma_{\mathrm{on}}$ ไปแล้ว (ตัวอย่างที่ $t \approx 21.49$ s:
$0.489 / 0.438 / 0.425 / 0.442$ ตามลำดับ IBR$_1$--IBR$_4$) จึงไม่มี dwell
สะสมใด ๆ up-stress กลับมาเสถียรอีกครั้งที่ $t = 21.99$ s (IBR$_2$ ขึ้นก่อน) และเมื่อครบ $T_{d,\mathrm{on}} = 0.10$ s ที่ $t = 22.089$ s
supervisor จึง commit การ augment ครั้งแรก

\subsection{ผลของครั้งแรก: เป้าหมายสองตัว สลับจริงตัวเดียว}
event log บันทึกไว้ชัดเจน:

\begin{quote}\small
\code{gfm\_support\_augment committed at t=22.089; selected=[2 4],
reference=2, 1 device(s) transitioned}
\end{quote}

\code{selected=[2 4]} คือ\emph{ชุดเป้าหมาย}ที่ตาราง (fingerprint
\code{f62cbf2f}) รับรองไว้: GFM สองตัวคือบัส 2 (IBR$_1$) และบัส 6
(IBR$_3$) แต่บันทึกว่า ``1 device(s) transitioned'' เพราะ IBR$_1$
\emph{เป็น GFM อยู่แล้ว} ตั้งแต่ $t = 20.0$ s --- event log แถว
\code{sg\_trip} บันทึก \code{selected=[2], reference=2, committed 1
authenticated selected-GFM resource(s)} หมายความว่าธุรกรรม trip ของ SG
ส่งมอบ angle reference ให้บัส 2 และสลับมันเป็น GFM ในการ commit เดียวกัน
(นี่คือแถบน้ำเงินยาว $t = 20.0$--$116.238$ s ของ IBR$_1$ ใน Figure 8)
ดังนั้นเมื่อ augment ที่ $t = 22.089$ s ต้องการชุด $\{2,4\}$ จึงเหลือ
ตัวที่ต้องสลับจริงเพียง device 4 (บัส 6, IBR$_3$) ---
\textbf{แถบแดงสั้น 22.089--25.303 s ใน Figure 8 คือ IBR$_3$ ไม่ใช่
IBR$_1$} และสังเกตว่า \emph{ผู้ทริกเกอร์} คือ IBR$_2$ (บัส 3)
แต่ \emph{ผู้ถูกเลือก} คือ IBR$_3$ (บัส 6) ตามกฎของตาราง
ไม่ใช่ตามค่า $S_i$

\section{ชั้นที่ 3: ทำไมจึงไม่มีการ augment ครั้งที่สอง}

หลัง $t = 22.089$ s ชุด GFM คือ $\{2,4\}$ (IBR$_1$ กับ IBR$_3$) ชุด GFL
จึงเหลือ $\{3,5\}$ คือ \textbf{IBR$_2$ (บัส 3) กับ IBR$_4$ (บัส 8)
เพียงสองตัว} และเงื่อนไข up-stress ถูกประเมินใหม่ \emph{เฉพาะบนชุดนี้}
พร้อมตัวจับเวลาที่ต้องเริ่มนับใหม่ทั้งหมด จากข้อมูลจริง:

\begin{center}\small
\begin{tabular}{@{}lccc@{}}
\toprule
$t$ [s] & $S$ IBR$_2$ (บัส 3) & $S$ IBR$_4$ (บัส 8) & up-stress บน GFL $\{3,5\}$ \\
\midrule
22.10 & 1.000 & 0.947 & จริง \\
22.51 & 0.678 & 0.606 & จริง \\
22.98 & 0.697 & 0.497 & จริง \\
23.08 & 0.668 & 0.477 & จริง \\
23.21 & 0.633 & 0.463 & \textbf{เท็จ} (ทุกตัวหลุด) \\
23.60 & 0.505 & 0.395 & เท็จ ต่อเนื่องจนสิ้น run \\
\bottomrule
\end{tabular}
\end{center}

หลัง commit มีช่วง up-stress เหลือสองช่วง: 22.089--22.219 s (129 ms) กับ
22.482--23.133 s (651 ms) --- ทว่า IBR$_4$ หลุดใต้ threshold ตั้งแต่
$t = 22.19$ s เหลือ IBR$_2$ ยืนอยู่ตัวเดียว และก่อนช่วงที่สองจะยาวถึง
$T_{d,\mathrm{on}} = 100$ ms จากการรีเซ็ต dwell หลัง commit (22.089 + 0.100
= 22.189 s) up-stress ก็ขาดช่วงไปแล้วที่ $t = 22.219$ s ทำให้ dwell รีเซ็ต
อีกครั้ง ส่วนช่วง 22.482--23.133 s แม้ยาว 651 ms แต่เกิดขึ้น\emph{หลัง}จาก
IBR$_3$ เป็น GFM แล้ว เกาะจึงฟื้นตัวเร็วขึ้นเรื่อย ๆ จน IBR$_2$ หลุดที่
$t = 23.133$ s ก่อน dwell ชุดใหม่จะยืนยันครบ --- หลัง $t = 23.2$ s
ทุกตัวต่ำกว่า 0.65 ถาวร supervisor เข้าสู่สถานะ hysteresis hold

เมื่อดัชนีทั้งระบบตกใต้ $\Gamma_{\mathrm{off}} = 0.35$ ครบทุกตัวที่
$t \approx 24.3$ s และค้างครบ $T_{d,\mathrm{off}} = 1.00$ s supervisor
จึง commit การ release ที่ $t = 25.303$ s ($\{2,4\} \to \{2\}$)
คืน IBR$_3$ กลับเป็น GFL --- นี่คือจุดสิ้นสุดของแถบแดงใน Figure 8
ส่วน IBR$_1$ (บัส 2) ถูกคงไว้เป็น GFM เพราะมันคือ angle-reference owner
(ผู้ถือ reference จะไม่ถูก release ตราบใดที่ SG ยังไม่กลับมา)

\section{ชั้นที่ 4: ช่วง fault $t = 60$ s --- spike ที่สูงที่สุดใน Figure 7}

ตอน fault bolted ที่บัส 9 ($Z_f = 0.01 + 0.01i$, 60.00--60.15 s) ดัชนีของ
ทั้งสี่ตัวถูกปักที่ $S_i = 1.000$ ตลอดฟอลต์ --- นี่คือแหล่งของเสาสูงสุดใน
Figure 7 แต่:

\begin{itemize}
\item ฟอลต์ยาวเพียง 150 ms และ $S_i$ ทุกตัวตกใต้ $\Gamma_{\mathrm{on}}$
      ที่ $t = 60.151$ s ทันทีหลังเคลียร์ ตัวจับเวลาได้ dwell ต่อเนื่อง
      $\approx 151$ ms (เกิน 100 ms อยู่หน่อยเดียว) ก่อนหลักฐานจะพับลง
\item ณ เวลานั้น IBR$_2$ กับ IBR$_4$ ยังเป็น GFL (IBR$_1$ ยังถือ
      reference เป็น GFM อยู่) ดังนั้นหาก dwell ครบก็จะเกิด augment
      รอบใหม่ --- แต่พอเคลียร์ฟอลต์ปุ๊บ ภายใน 10 ms ค่า $S_i$
      ร่วงเหลือ 0.06--0.17 และต่อไปต่ำกว่า 0.27 ตลอด เกาะพิสูจน์ตัวเองว่า
      \emph{แข็งแรงโดยไม่ต้องเพิ่ม former เลย}
\item ตรรกะจึงตอบ ``ไม่ทำอะไร'' --- ซึ่งคือพฤติกรรมที่ถูกต้อง:
      severity index สูงชั่วคราวจากฟอลต์สั้นไม่ใช่เหตุให้เปลี่ยน
      topology ของ forming
\end{itemize}

ส่วนจุดสิ้นสุดของแถบ GFM ของ IBR$_1$ ใน Figure 8 ($t = 116.238$ s)
เกิดจากเส้นทางแยกของ SG-reclose: เมื่อ SG กลับมาที่ $t = 102.363$ s
reference ถูกย้ายคืนให้ SG ทันที แต่ IBR$_1$ ยังคงโหมด GFM ไว้อีก
$\approx 13.9$ s จนกว่าดัชนีของมันจะคงอยู่ใต้ $\Gamma_{\mathrm{off}}$ ครบ
เกณฑ์ supervisor จึงคืนมันเป็น GFL ผ่าน \code{sg\_reselection}
ทิ้งเกาะให้เป็น GFL ทั้งหมดจนจบ run

\section{สรุป}

\begin{enumerate}[label=\textbf{\arabic*.}]
\item กราฟ $S_i$ ที่ทะลุเส้น $\Gamma_{\mathrm{on}}$ ใน Figure 7 คือ
      \emph{หลักฐานนำเข้า} ของการตัดสินใจระดับระบบ ไม่ใช่คำสั่งรายตัว
\item IBR$_1$ ไม่ได้ถูก selector เลือกจากค่า $S_i$ --- มันถูกสลับโดย
      \emph{ธุรกรรม trip ของ SG เอง} ตอน $t = 20$ s เพื่อรับ angle
      reference (แถบน้ำเงิน) ส่วนการ augment ที่ $t = 22.089$ s เลือกชุด
      $\{2,4\}$ จากตาราง ซึ่งหมายถึงตัวที่สลับจริงคือ IBR$_3$
      (แถบแดง) เพราะ IBR$_1$ เป็น GFM อยู่แล้ว
\item IBR$_2$ กับ IBR$_4$ ไม่ได้ ``โดนขัด'' --- พวกมันแพ้เกตเวลา:
      ช่วงแรกโดน minimum-hold 1 s กลืนหายไป ช่วงที่เหลือ (22.10--23.13 s)
      dwell ต่อเนื่องไม่เคยถึง 100 ms หลังจาก IBR$_3$ ถูกเสริมเข้าไปแล้ว
      ดัชนีจึงสลายตัวก่อนเกตจะเปิดอีก
\item พฤติกรรมทั้งหมดสอดคล้องกับหลักออกแบบ:
      \emph{สลับให้น้อยที่สุดเท่าที่จำเป็น และพิสูจน์ก่อนทุกครั้ง}
      --- Figure 7 กับ Figure 8 จึงไม่ขัดกัน แต่เป็นภาพสองมุมของ
      กระบวนการเดียวกัน
\end{enumerate}

\end{document}
